\bookchapter{Capitolo 57: La strada del livello successivo — permessi, notifiche e architettura moderna}

\chapterintro{Saper costruire un'app è solo l'inizio. Questo capitolo affronta gli ostacoli che portano da «saper fare» a «fare bene»: come chiedere i permessi con eleganza, come inviare notifiche senza dare fastidio, come alzare di un livello l'architettura e — nel 2026 — come scrivere codice Android in coppia con l'AI. Ogni punto ti capiterà subito fra le mani in un progetto vero.}

\booksection{Permessi a runtime: prima bussi, poi entri}

Fotocamera, posizione, contatti, notifiche… queste «zone sensibili» Android vuole che l'app le chieda all'utente a runtime (da Android 6 in poi). La regola: prima dichiari (nel Manifest), poi richiedi (nel codice) e, se l'utente rifiuta, degradi con eleganza. In Compose la via più tranquilla è rememberPermissionState della libreria ufficiale Accompanist:

\begin{lstlisting}[language=Kotlin]
// Prima dichiaralo nel Manifest:
// <uses-permission android:name="android.permission.CAMERA" />

@Composable
fun CameraButton() {
    val permission = rememberPermissionState(
        permission = android.Manifest.permission.CAMERA
    )
    Button(onClick = {
        when {
            permission.status.isGranted -> openCamera()
            permission.status.shouldShowRationale -> showExplainDialog()
            else -> permission.launchPermissionRequest()
        }
    }) { Text("Scatta") }
}
\end{lstlisting}

\begin{tcbwarn}{Attenzione: e se il permesso viene rifiutato per sempre?}
Dopo che l'utente ha spuntato «non chiedere più», launchPermissionRequest non apre più il popup. A quel punto guida l'utente ad attivarlo a mano nelle impostazioni di sistema (con un Intent che apre la pagina delle impostazioni dell'app) e spiega nell'interfaccia «a cosa serve questo permesso». Con un atteggiamento gentile la voglia di concedere il permesso cresce parecchio: il testo con cui chiedi i permessi merita di essere scritto sul serio.

\end{tcbwarn}

\booksection{Notifiche: dire cose utili al momento giusto}

Le notifiche sono il canale con cui l'app raggiunge l'utente, e anche la funzione che gli utenti mettono più volentieri al bando. Tre regole ferree: manda solo cose utili, lascia all'utente la possibilità di spegnerle, non disturbare a notte fonda. Uso di base:

\begin{lstlisting}[language=Kotlin]
// Su Android 13+ devi prima richiedere il permesso POST_NOTIFICATIONS (vedi la sezione precedente)

// 1. Crea il canale di notifica (obbligatorio da Android 8+)
val channel = NotificationChannel(
    "reminder", "Promemoria", NotificationManager.IMPORTANCE_DEFAULT
)
val manager = getSystemService(NotificationManager::class.java)
manager.createNotificationChannel(channel)

// 2. Invia la notifica
val notification = NotificationCompat.Builder(this, "reminder")
    .setSmallIcon(R.drawable.ic_notification)
    .setContentTitle("Promemoria spese")
    .setContentText("Oggi non hai ancora registrato spese")
    .setPriority(NotificationCompat.PRIORITY_DEFAULT)
    .build()
if (ActivityCompat.checkSelfPermission(this,
        Manifest.permission.POST_NOTIFICATIONS) == PackageManager.PERMISSION_GRANTED
) {
    manager.notify(1, notification)
}
\end{lstlisting}

Le attività programmate (per esempio «ricorda di segnare le spese ogni giorno alle 20:00») si affidano a WorkManager: è il «maggiordomo di pianificazione in background» raccomandato ufficialmente, garantisce che i compiti partano al momento giusto e ti libera dalle preoccupazioni sull'ottimizzazione della batteria. Per esigenze tipo sveglia non sbagli.

\booksection{Architettura avanzata: dependency injection e modularizzazione}

Quando il progetto cresce, si presentano due «dolci seccature»:

\begin{itemize}
  \item Creazione di oggetti in disordine: ogni schermata si costruisce da sola database e API — con Hilt (framework di dependency injection) si amministra tutto in un posto solo: «chiunque chieda qualcosa, qualcuno glielo porge»;
  \item Codice che si accumula a montagna: un progetto a modulo unico si modifica sempre più lentamente — dividilo in più moduli (app / core / feature-xxx), con ogni funzione sviluppata e testata per conto suo.
\end{itemize}

\begin{tcbinfo}{Consiglio: non progettare troppo}
Hilt e la modularizzazione sono ottimi strumenti, ma in un progetto in solitaria non avere fretta. Per il registro delle spese bastano e avanzano un modulo unico con dipendenze create a mano. Introducili quando il codice supera le decine di migliaia di righe, la compilazione comincia a rallentare o ci si mette più persone. Gli strumenti servono all'obiettivo, non al fare sfoggio — è anche il principio «semplicità sopra ogni cosa» del nostro libro sulla programmazione con l'AI.

\end{tcbinfo}

\booksection{La novità del 2026: sviluppare Android con l'assistenza dell'AI}

Nel 2026 l'AI è ormai incastonata nel flusso di lavoro Android. Se vuoi accelerare, questa combinazione è ben rodata:

\begin{center}
\begin{tabularx}{\linewidth}{|X|X|X|}
\hline
\textbf{Fase} & \textbf{Come aiuta l'AI} & \textbf{Il tuo compito} \tabularnewline
\hline
Scrivere codice & Completamento e generazione con Gemini integrato in Android Studio o con plugin AI di terze parti & Descrivere bene i requisiti, revisionare i risultati \tabularnewline
Cercare informazioni & Consegnare all'AI documentazione ufficiale e messaggi di errore da interpretare & Verificare le conclusioni chiave sulla documentazione ufficiale \tabularnewline
Scrivere test & Fare generare all'AI i casi di test nelle tre categorie «normale/limite/anomalo» & Confermare che i casi rispettino il senso del dominio \tabularnewline
Correggere i bug & Incollare stacktrace e codice correlato: l'AI localizza e ripara & Capire le modifiche e verificare facendo girare i test \tabularnewline
\hline
\end{tabularx}
\end{center}

\begin{center}\small\color{gray}I quattro scenari dello sviluppo Android assistito dall'AI\end{center}

La metodologia concreta — come dare istruzioni all'AI, come accettare i risultati, come schivare le trappole — occupa ben 11 capitoli nella parte 7 di questo libro, «Programmazione AI in pratica» (capitoli 37-47): è fortemente consigliato leggerli in parallelo. La sintassi Android la copre questa parte, i segreti della collaborazione con l'AI la parte 7.

\booksection{Ottimizzazione delle prestazioni di Compose e strumenti di debug}

Con Compose si scrive volentieri, ma troppa «ricomposizione» può far inciampare. Tre trucchetti pratici:

\begin{itemize}
  \item Parametri stabili: aggiungi @Immutable alle data class o usa derivedStateOf, per evitare che variazioni di stato irrilevanti scatenino ricomposizioni su larga scala;
  \item Ricorda di usare la key: items(list, key = \{ it.id \}) fa sì che aggiunte e cancellazioni tocchino solo le voci interessate — senza key, cancellare una voce può ridisegnare l'intera lista;
  \item La modalità Compose del Layout Inspector: mostra in tempo reale quante volte viene ricomposto ogni Composable; chi si ricompiona troppo si vede a colpo d'occhio e ottimizzerai puntando sui punti caldi.
\end{itemize}

In più, apri i report del compilatore Compose (Metrics): generano un resoconto di «quali componenti non sono stabili» — sulla via del livello avanzato, i dati sono sempre più affidabili delle sensazioni.

\begin{tcbpractice}{Mani in pasta: aggiungi il «tris» al registro delle spese}
(1) Aggiungi al registro delle spese il promemoria «segnare le spese ogni giorno alle 20:00» (WorkManager + canale di notifica + permesso notifiche di Android 13); (2) passa la creazione di database e API all'iniezione con Hilt; (3) prova a usare uno strumento AI per realizzare più in fretta la funzione «filtro mensile», poi controlla il risultato col metodo di collaudo a tre livelli. Fatto il tris, il tuo progetto avrà un aspetto decisamente professionale.

\end{tcbpractice}

\begin{tcbsummary}{Riassunto del capitolo}
\begin{itemize}
  \item Permessi sensibili: dichiarazione nel Manifest + richiesta a runtime + guida elegante dopo il rifiuto
  \item Le tre regole delle notifiche: solo cose utili, spegnibili, nessun fastidio; canale + permesso, non possono mancare
  \item Le attività programmate con WorkManager, non reinventare la ruota
  \item Hilt amministra gli oggetti, la modularizzazione il codice — nei progetti piccoli in solitaria non avere fretta
  \item Quattro scenari AI: scrivere codice, cercare informazioni, scrivere test, correggere bug; la metodologia è nella parte 7 (capitoli 37-47)
\end{itemize}

\end{tcbsummary}
